The transport principle completes $\Cl(p,q)$ to $\Cl(p{+}1,q{+}1)$ by adjoining one hyperbolic pair, and the completed algebra carries a null cone of its own, so the completion iterates: the $k$-th rung is $T_{k} = \Cl(p{+}k, q{+}k)$, with adjoined null pairs $v_{j}, w_{j}$ satisfying, as vectors, $v_{i}v_{j} + v_{j}v_{i} = 0$, $w_{i}w_{j} + w_{j}w_{i} = 0$ and $v_{i}w_{j} + w_{j}v_{i} = \delta_{ij}$. Inside $T_{k}$ the subalgebra $B_{k}$ generated by the $2k$ plane generators is $\Cl(k,k) \cong \mathrm{M}_{2^{k}}(\R)$, the twisted ground generators generate a commuting copy $A_{k} \cong \Cl(p,q)$, and multiplication $A_{k} \otimes B_{k} \to T_{k}$ is an isomorphism, $T_{k} \cong \Cl(p,q) \otimes \mathrm{M}_{2^{k}}(\R)$. With $W_{k} = \operatorname{span}\{v_{1}, \ldots, v_{k}\}$ the Chevalley action makes $\Lambda^{\bullet}W_{k}$, of dimension $2^{k}$, an irreducible $B_{k}$-module, the vacuum $1$ annihilated by every $w_{j}$ and cyclic for the $v_{j}$, and each rung doubles it: $\Lambda^{\bullet}W_{k+1} = \Lambda^{\bullet}W_{k} \otimes \Lambda^{\bullet}(\R v_{k+1})$, one qubit per rung. The Witt index rises by exactly one per rung, and the Witt index of a real quadratic space is the maximal number of anticommuting mode pairs realisable by vectors. On each complexified floor there is a unique conjugate-linear antiautomorphism $*$ fixing the positive frame generators and negating the negative ones, under which the ladder pairs are mutually adjoint, $v_{j}^{*} = w_{j}$. \emph{Provenance:} \cite[Def.~2.2, Thm.~3.2, Prop.~4.1, Thms.~4.2, 4.3, Lem.~5.4]{ButlerT3} (\claimref{tower.structure-tower-c-norms}, \claimref{tower.witt-index-quantum-capacity}, \claimref{tower.commuting-copies}, \claimref{tower.fock-module-one-qubit-per-rung}, \claimref{tower.ladder-pairs-satisfy-car-vectors}).
